\clearpage
\section{TX-side Conversion Gain --- Lower Sideband (Spectral Flip)}
\label{sec:tx-gain-lower}

\textbf{What is conversion gain?} Exactly as in Section~\ref{sec:tx-gain-upper}: the power of the signal coming out of the ``TX-side FR3 RF Out'' connector, minus the power of the signal fed into the ``TX-side FR1 IF In'' connector, both in dBm, gives
the conversion gain in dB. Only the sideband changes here, now the lower one instead of the upper.

\textbf{What's plotted.} Identical methodology to Section~\ref{sec:tx-gain-upper}, with the unit tuned to the
lower sideband instead of the upper: each of the first seven figures holds FR1 fixed at one center value (1.0
up to the unit's 4.0\,GHz maximum, in 0.5\,GHz steps) and plots conversion gain against the full measured FR3
range for that input.\footnote{Every point on both traces is a real measurement, never a simulation or an
interpolation between sparse samples: the FieldFox's built-in tracking generator drives FR1 while the
FieldFox's own receiver tracks the corresponding FR3 output in lock-step, at a calibrated −40\,dBm input
level.} Grey is the raw amplifier-and-mixer chain with the AM3153 filter bypassed; blue is the filter engaged.
Both traces plot the \emph{average} conversion gain across each point's own real swept window, nominally
500\,MHz wide (the same for every point on every panel) --- mean linear power, converted back to dB, not a
single CW tone: what a real wideband signal occupying that window would experience. The gap between the two
traces, at any point, is exactly how much the filter costs in insertion loss there, same as in
Section~\ref{sec:tx-gain-upper}.

\textbf{Measurement conditions.} \lmxDriveRule{} (set from the LO frequency), with the LO leakage suppressed by the unit's stored LO-null calibration (Section~\ref{sec:lo-suppression}), exactly as in normal use. The unit was measured hot, as the worst case: no fan, LMX2820 die at \condGainMinC--\condGainMaxC\textdegree C. A cooler unit will typically do somewhat better.

\textbf{How the filter is chosen.} The same two-step rule as Section~\ref{sec:tx-gain-upper}: first find
whichever of the 48 candidate configurations (3 bands × 16 sub-states) passes the most desired-sideband
energy on its own, ignoring the image; then, restricted to candidates within 3\,dB of that best desired power,
choose whichever rejects the image the most. Consecutive points that land on the same filter form a
\emph{region}; the dashed vertical lines on each panel mark every boundary between regions, and the blue hex
code running up from the x-axis inside each region (e.g.\ \texttt{0x30}) is the actual \texttt{filter\_bits}
value used for every point in it --- so, just as in Section~\ref{sec:tx-gain-upper}, the exact filter behind
any measurement on these panels is always readable directly off the figure, with no need to consult a separate
table.

\textbf{Window width.} Every point's window is exactly 500\,MHz wide, the same fixed window used throughout
Section~\ref{sec:tx-gain-upper}: a fixed filter setting's measured response across that full window is exactly
what a real wideband signal at that setting would experience.\footnote{As in Section~\ref{sec:tx-gain-upper},
Figure \ref{fig:2-7} (FR1 = 4.0\,GHz) is the one exception: its window is 3.75--4.25\,GHz, and 250\,MHz of that
exceeds the unit's 0.7--4.0\,GHz IF design spec --- called out directly on that panel and outlined in red on
Figures \ref{fig:2-8} and \ref{fig:2-9}.} For the lower sideband, the masked (invalid-LO) region of the heatmaps
sits at the opposite end of the FR3 axis from the upper sideband, since LO = FR3 + FR1 here rather than
FR3 − FR1. Those two heatmaps repeat the bypassed and optimal data across the full measured (FR1, FR3) grid,
for a single-glance view instead of the seven separate slices above, exactly as in Section~\ref{sec:tx-gain-upper}.

\figpair{images/sec2_fig1.pdf}{FR1 = 1.0\,GHz.}{fig:2-1}%
        {images/sec2_fig2.pdf}{FR1 = 1.5\,GHz.}{fig:2-2}
\figpair{images/sec2_fig3.pdf}{FR1 = 2.0\,GHz.}{fig:2-3}%
        {images/sec2_fig4.pdf}{FR1 = 2.5\,GHz.}{fig:2-4}
\figpair{images/sec2_fig5.pdf}{FR1 = 3.0\,GHz.}{fig:2-5}%
        {images/sec2_fig6.pdf}{FR1 = 3.5\,GHz.}{fig:2-6}

\figonesmall{images/sec2_fig7.pdf}{FR1 = 4.0\,GHz. 250\,MHz of this window falls outside the unit's IF design spec (shown in red).}{fig:2-7}
\figrow{images/sec2_fig8.pdf}{Conversion gain, lower sideband: full measured (FR1, FR3) grid, filter bypassed.
Zoom in to read each cell's exact values, printed as (FR1, FR3, conversion gain), with FR1 and FR3 in GHz and
the conversion gain in dB.}{fig:2-8}
\figrow{images/sec2_fig9.pdf}{Conversion gain, lower sideband: full measured (FR1, FR3) grid, optimal filter.
Zoom in to read each cell's exact values, printed as (FR1, FR3, conversion gain), with FR1 and FR3 in GHz and
the conversion gain in dB.}{fig:2-9}
